\chapter{Eight-Vertex Model and the XYZ Chain}
\label{chap:eight-vertex-xyz}

The six-vertex model is the vertex-model counterpart of the XXZ chain.  Its ice
rule gives a local $U(1)$ conservation law, and the trigonometric $R$-matrix gives
an interacting but integrable quantum spin chain.  The eight-vertex model is the
next member in the same hierarchy.  It relaxes the ice rule while preserving an
even-arrow, or $\ZZ_2$, constraint.  Algebraically this replaces the six-vertex
$R$-matrix by an elliptic $R$-matrix.  The Hamiltonian limit is no longer XXZ but
the fully anisotropic XYZ spin chain.

The exact correspondence is
\[
\begin{aligned}
    \text{eight-vertex transfer matrix}
    &\quad\Longrightarrow\quad
    \text{elliptic Yang--Baxter equation}\\
    &\quad\Longrightarrow\quad
    \text{XYZ Hamiltonian}.
\end{aligned}
\]
The important conceptual change from the previous chapter is that the model is
still Yang--Baxter integrable, but it is no longer organized by total
magnetization sectors.  The missing $U(1)$ symmetry is the main reason why the
solution of the eight-vertex model is technically harder than the six-vertex
case.

This chapter does not derive Baxter's full solution.  Its goal is to set up the
model, state the elliptic $R$-matrix structure, and derive the Hamiltonian limit
that produces the XYZ chain.  Baxter's $TQ$ relation and the full spectral theory
belong to a more advanced treatment of elliptic integrability.

\section{Eight allowed vertices and local weights}
\label{sec:eight-vertex-local-weights}

The eight-vertex model is a classical statistical model on the square lattice.
As in the six-vertex model, arrows live on edges and a Boltzmann weight is
assigned to each vertex.  The local constraint is weaker than the ice rule:
\begin{equation}
\label{eq:eight-vertex-even-arrow-rule}
    \text{an allowed vertex has an even number of arrows pointing in.}
\end{equation}
Equivalently, the allowed numbers of incoming arrows are $0$, $2$, and $4$.
There are therefore eight allowed local configurations.  The six ice-rule
vertices are included, and the two additional vertices have all arrows pointing
in or all arrows pointing out.

In spin notation a local vertex weight is an operator on
$V\otimes V$, where $V\simeq\CC^2$.  We use the same ordered basis as in the
six-vertex chapter,
\begin{equation}
    \ket{\uparrow\uparrow},\quad
    \ket{\uparrow\downarrow},\quad
    \ket{\downarrow\uparrow},\quad
    \ket{\downarrow\downarrow}.
\end{equation}
A standard symmetric eight-vertex weight matrix has the form
\begin{equation}
\label{eq:eight-vertex-R-abcd}
    R=
    \begin{pmatrix}
        a&0&0&d\\
        0&b&c&0\\
        0&c&b&0\\
        d&0&0&a
    \end{pmatrix}.
\end{equation}
The six-vertex model is recovered when $d=0$.  The extra $d$-terms are the new
eight-vertex processes
\begin{equation}
\label{eq:eight-vertex-pair-flip-processes}
    \ket{\uparrow\uparrow}\longleftrightarrow\ket{\downarrow\downarrow}.
\end{equation}
They change the total $S^z$ of the two lines by two units.  Thus the model no
longer conserves total $S^z$, but it conserves its parity.

\begin{proposition}[$\ZZ_2$ conservation]
Let
\begin{equation}
    Q_2=\sigma^z\otimes\sigma^z .
\end{equation}
For the eight-vertex weight \cref{eq:eight-vertex-R-abcd},
\begin{equation}
\label{eq:eight-vertex-local-Z2-conservation}
    [R,Q_2]=0 .
\end{equation}
However, for $d\neq0$,
\begin{equation}
    \left[R,\sigma^z\otimes\id+\id\otimes\sigma^z\right]\neq0 .
\end{equation}
Thus the local $U(1)$ conservation of the six-vertex model is reduced to a
$\ZZ_2$ parity conservation.
\end{proposition}

\begin{proof}
The diagonal $a$ and $b$ processes clearly preserve both the total $z$-spin and
its parity.  The $c$ processes exchange
$\ket{\uparrow\downarrow}$ and $\ket{\downarrow\uparrow}$, so they also preserve
total $z$-spin.  The $d$ processes exchange
$\ket{\uparrow\uparrow}$ and $\ket{\downarrow\downarrow}$, changing total $S^z$
by $\pm2$ but preserving the eigenvalue of $\sigma^z\otimes\sigma^z$.  This is
exactly \cref{eq:eight-vertex-local-Z2-conservation}.
\end{proof}

\begin{remark}[Why the eight-vertex model is harder]
The six-vertex transfer matrix decomposes into fixed-magnetization sectors.  This
sector decomposition is the basic input behind the usual algebraic Bethe ansatz
with a simple ferromagnetic reference state.  In the eight-vertex model the
$d$-processes destroy $S^z$ conservation, so the standard highest-weight
construction no longer applies directly.  Baxter's solution replaces it with an
elliptic functional-relation method.
\end{remark}

\section{Row-to-row transfer matrix}
\label{sec:eight-vertex-transfer-matrix}

The transfer-matrix construction is formally the same as in the six-vertex
model.  On a row of length $L$, the row space is
\begin{equation}
\label{eq:eight-vertex-row-space}
    \calH_L=V_1\otimes V_2\otimes\cdots\otimes V_L,
    \qquad V_j\simeq\CC^2 .
\end{equation}
Introduce an auxiliary space $V_a\simeq\CC^2$.  Given a spectral-parameter family
$R(u)$ of eight-vertex weights, define the monodromy matrix
\begin{equation}
\label{eq:eight-vertex-monodromy}
    M_a(u)=R_{aL}(u)R_{a,L-1}(u)\cdots R_{a1}(u),
\end{equation}
and the row-to-row transfer matrix
\begin{equation}
\label{eq:eight-vertex-transfer-matrix}
    \tau(u)=\Tr_a M_a(u)
    =\Tr_a\!\left[R_{aL}(u)R_{a,L-1}(u)\cdots R_{a1}(u)\right] .
\end{equation}
For an $L\times M$ torus,
\begin{equation}
\label{eq:eight-vertex-partition-function}
    Z_{L,M}(u)=\Tr_{\calH_L}\tau(u)^M .
\end{equation}
The spectrum of $\tau(u)$ controls the thermodynamic limit of the classical
model.

Because of the local parity conservation in
\cref{eq:eight-vertex-local-Z2-conservation}, the transfer matrix commutes with
the global spin-parity operator
\begin{equation}
\label{eq:eight-vertex-global-parity}
    \mathcal{P}=\prod_{j=1}^{L}\sigma_j^z,
    \qquad
    [\tau(u),\mathcal{P}]=0 .
\end{equation}
For $d\neq0$ one does not generally have
\begin{equation}
    [\tau(u),S^z_{\rm tot}]=0,
    \qquad
    S^z_{\rm tot}=\frac12\sum_{j=1}^{L}\sigma_j^z .
\end{equation}
This is the transfer-matrix reflection of the fact that the corresponding XYZ
Hamiltonian has only a discrete spin-flip/parity symmetry, not a continuous
$U(1)$ symmetry.

\section{Baxter's elliptic parametrization}
\label{sec:baxter-elliptic-parametrization}

A generic matrix of the form \cref{eq:eight-vertex-R-abcd} does not automatically
produce commuting transfer matrices.  Integrability requires that the weights be
embedded in an elliptic spectral-parameter family satisfying the Yang--Baxter
equation.

Let $\vartheta_1(u)$ and $\vartheta_4(u)$ denote Jacobi theta functions with a
fixed nome $q$, $\abs{q}<1$.  Let $\eta$ be the crossing parameter.  One common
Baxter parametrization, shifted so that the regular point is at $u=0$, is
\begin{equation}
\label{eq:eight-vertex-elliptic-weights}
\begin{aligned}
    a(u)&=\rho\,\vartheta_4(2\eta)\,\vartheta_4(u)\,\vartheta_1(u+2\eta),\\
    b(u)&=\rho\,\vartheta_4(2\eta)\,\vartheta_1(u)\,\vartheta_4(u+2\eta),\\
    c(u)&=\rho\,\vartheta_1(2\eta)\,\vartheta_4(u)\,\vartheta_4(u+2\eta),\\
    d(u)&=\rho\,\vartheta_1(2\eta)\,\vartheta_1(u)\,\vartheta_1(u+2\eta).
\end{aligned}
\end{equation}
Here $\rho$ is an arbitrary nonzero scalar normalization.  The corresponding
$R$-matrix is
\begin{equation}
\label{eq:eight-vertex-R-u}
    R(u)=
    \begin{pmatrix}
        a(u)&0&0&d(u)\\
        0&b(u)&c(u)&0\\
        0&c(u)&b(u)&0\\
        d(u)&0&0&a(u)
    \end{pmatrix}.
\end{equation}
Since $\vartheta_1(0)=0$ and $\vartheta_4'(0)=0$, the regular point satisfies
\begin{equation}
\label{eq:eight-vertex-regularity}
    b(0)=d(0)=0,
    \qquad
    a(0)=c(0)=\rho_0,
\end{equation}
where
\begin{equation}
\label{eq:eight-vertex-rho0}
    \rho_0
    =\rho\,\vartheta_4(0)\vartheta_4(2\eta)\vartheta_1(2\eta).
\end{equation}
Therefore
\begin{equation}
\label{eq:eight-vertex-R0-P}
    R(0)=\rho_0 P,
\end{equation}
where $P$ is the two-spin permutation operator,
\begin{equation}
\label{eq:eight-vertex-permutation-operator}
    P\ket{\alpha}\otimes\ket{\beta}
    =\ket{\beta}\otimes\ket{\alpha}.
\end{equation}

\begin{remark}[Normalization and shift conventions]
The eight-vertex literature contains several equivalent theta-function
conventions.  Some place the regular point at $u=\eta$ rather than $u=0$, and
some use theta functions denoted by $H$ and $\Theta$.  The convention in
\cref{eq:eight-vertex-elliptic-weights} is chosen to match the Hamiltonian-limit
convention used in these notes: $R(0)$ is proportional to the permutation
operator.
\end{remark}

For the family \cref{eq:eight-vertex-R-u}, the Yang--Baxter equation takes the
same formal form as in the six-vertex case:
\begin{equation}
\label{eq:eight-vertex-YBE}
    R_{12}(u-v)R_{13}(u-w)R_{23}(v-w)
    =
    R_{23}(v-w)R_{13}(u-w)R_{12}(u-v).
\end{equation}
This identity is a nontrivial theta-function identity.  Its consequence is the
commutativity of transfer matrices:
\begin{equation}
\label{eq:eight-vertex-commuting-transfer-matrices}
    [\tau(u),\tau(v)]=0
    \qquad
    \text{for all }u,v .
\end{equation}
Thus the eight-vertex model is integrable in the same structural sense as the
six-vertex model: there is a one-parameter commuting family of transfer
matrices.

\begin{principle}[Elliptic integrability]
The hierarchy
\begin{equation}
    \text{rational }R\text{-matrix}
    \longrightarrow
    \text{trigonometric }R\text{-matrix}
    \longrightarrow
    \text{elliptic }R\text{-matrix}
\end{equation}
corresponds, respectively, to the XXX, XXZ, and XYZ spin chains.  The eight-vertex
model is the elliptic vertex model whose Hamiltonian limit is the XYZ chain.
\end{principle}

A useful way to see the difference from the six-vertex model is the Pauli-basis
form.  Since
\begin{align}
    \id\otimes\id
    &\mapsto
    \begin{pmatrix}
        1&0&0&0\\0&1&0&0\\0&0&1&0\\0&0&0&1
    \end{pmatrix},\\
    \sigma^x\otimes\sigma^x
    &\mapsto
    \begin{pmatrix}
        0&0&0&1\\0&0&1&0\\0&1&0&0\\1&0&0&0
    \end{pmatrix},\\
    \sigma^y\otimes\sigma^y
    &\mapsto
    \begin{pmatrix}
        0&0&0&-1\\0&0&1&0\\0&1&0&0\\-1&0&0&0
    \end{pmatrix},\\
    \sigma^z\otimes\sigma^z
    &\mapsto
    \begin{pmatrix}
        1&0&0&0\\0&-1&0&0\\0&0&-1&0\\0&0&0&1
    \end{pmatrix},
\end{align}
we may rewrite \cref{eq:eight-vertex-R-u} as
\begin{equation}
\label{eq:eight-vertex-R-Pauli-basis}
    R(u)=
    A(u)\,\id\otimes\id
    +B(u)\,\sigma^x\otimes\sigma^x
    +C(u)\,\sigma^y\otimes\sigma^y
    +D(u)\,\sigma^z\otimes\sigma^z,
\end{equation}
with
\begin{equation}
\label{eq:eight-vertex-ABCD}
    A(u)=\frac{a(u)+b(u)}{2},
    \qquad
    D(u)=\frac{a(u)-b(u)}{2},
\end{equation}
\begin{equation}
\label{eq:eight-vertex-BCD}
    B(u)=\frac{c(u)+d(u)}{2},
    \qquad
    C(u)=\frac{c(u)-d(u)}{2}.
\end{equation}
The distinction between $\sigma^x\sigma^x$ and $\sigma^y\sigma^y$ is controlled
by the $d$-weight.  When $d=0$, the coefficients of
$\sigma^x\otimes\sigma^x$ and $\sigma^y\otimes\sigma^y$ coincide, and the model
reduces to the six-vertex/XXZ type.

\section{Hamiltonian limit and the XYZ chain}
\label{sec:eight-vertex-xyz-hamiltonian-limit}

The Hamiltonian limit is obtained from the logarithmic derivative of the transfer
matrix at the regular point.  Because $R(0)=\rho_0P$, the transfer matrix at
$u=0$ is proportional to the one-site translation operator.  The first derivative
therefore gives a sum of local nearest-neighbour terms.

\begin{proposition}[Hamiltonian from the eight-vertex transfer matrix]
Let $\tau(u)$ be the periodic transfer matrix built from the regular
$R$-matrix \cref{eq:eight-vertex-R-u}.  Define
\begin{equation}
\label{eq:eight-vertex-H-log-derivative}
    H_{\rm XYZ}
    =\left.\frac{\dd}{\dd u}\log\tau(u)\right|_{u=0}.
\end{equation}
After dropping an additive scalar multiple of the identity and allowing an
overall energy-scale normalization, one obtains
\begin{equation}
\label{eq:XYZ-Hamiltonian-main}
    H_{\rm XYZ}
    =
    \sum_{j=1}^{L}
    \left(
        J_x\,\sigma_j^x\sigma_{j+1}^x
        +J_y\,\sigma_j^y\sigma_{j+1}^y
        +J_z\,\sigma_j^z\sigma_{j+1}^z
    \right),
    \qquad
    \sigma_{L+1}^{\alpha}\equiv\sigma_1^{\alpha}.
\end{equation}
The three couplings are determined by the first derivatives of the vertex
weights at the regular point:
\begin{equation}
\label{eq:eight-vertex-Jxyz-from-weight-derivatives}
    J_x=\frac{b'(0)+d'(0)}{2\rho_0},
    \qquad
    J_y=\frac{b'(0)-d'(0)}{2\rho_0},
    \qquad
    J_z=\frac{a'(0)-c'(0)}{2\rho_0}.
\end{equation}
The discarded additive constant is
\begin{equation}
\label{eq:eight-vertex-additive-constant}
    E_0=\frac{a'(0)+c'(0)}{2\rho_0}
\end{equation}
per bond.
\end{proposition}

\begin{proof}
By regularity, $R(0)=\rho_0P$.  The standard transfer-matrix argument gives a
local logarithmic derivative
\begin{equation}
\label{eq:eight-vertex-local-h-P-Rprime}
    h_{j,j+1}=\frac{1}{\rho_0}P_{j,j+1}R'_{j,j+1}(0).
\end{equation}
For the matrix form \cref{eq:eight-vertex-R-u},
\begin{equation}
\label{eq:eight-vertex-PRprime-matrix}
    \frac{1}{\rho_0}PR'(0)
    =\frac{1}{\rho_0}
    \begin{pmatrix}
        a'(0)&0&0&d'(0)\\
        0&c'(0)&b'(0)&0\\
        0&b'(0)&c'(0)&0\\
        d'(0)&0&0&a'(0)
    \end{pmatrix}.
\end{equation}
Comparing this matrix with
\begin{equation}
\label{eq:eight-vertex-local-XYZ-comparison}
    E_0\id
    +J_x\sigma^x\otimes\sigma^x
    +J_y\sigma^y\otimes\sigma^y
    +J_z\sigma^z\otimes\sigma^z
\end{equation}
gives exactly \cref{eq:eight-vertex-Jxyz-from-weight-derivatives} and
\cref{eq:eight-vertex-additive-constant}.
\end{proof}

For the theta-function weights in \cref{eq:eight-vertex-elliptic-weights}, these
formulas give
\begin{equation}
\label{eq:eight-vertex-JxJy-theta}
    J_x
    =
    \frac{\vartheta_1'(0)}{2\vartheta_4(0)\vartheta_4(2\eta)\vartheta_1(2\eta)}
    \left[\vartheta_4(2\eta)^2+\vartheta_1(2\eta)^2\right],
\end{equation}
\begin{equation}
\label{eq:eight-vertex-Jy-theta}
    J_y
    =
    \frac{\vartheta_1'(0)}{2\vartheta_4(0)\vartheta_4(2\eta)\vartheta_1(2\eta)}
    \left[\vartheta_4(2\eta)^2-\vartheta_1(2\eta)^2\right],
\end{equation}
and
\begin{equation}
\label{eq:eight-vertex-Jz-theta}
    J_z
    =
    \frac12
    \left[
        \frac{\vartheta_1'(2\eta)}{\vartheta_1(2\eta)}
        -
        \frac{\vartheta_4'(2\eta)}{\vartheta_4(2\eta)}
    \right].
\end{equation}
These expressions depend on the chosen normalization of the spectral parameter.
A rescaling $u\mapsto\lambda u$ rescales all three couplings by the common factor
$\lambda$.  The physically meaningful anisotropies are the coupling ratios.

\begin{remark}[Sign convention]
Some condensed-matter texts write
\begin{equation}
    H=-\sum_j
    \left(
        J_x\sigma_j^x\sigma_{j+1}^x
        +J_y\sigma_j^y\sigma_{j+1}^y
        +J_z\sigma_j^z\sigma_{j+1}^z
    \right).
\end{equation}
This differs from \cref{eq:XYZ-Hamiltonian-main} by redefining
$J_{\alpha}\mapsto -J_{\alpha}$.  The integrable structure is unaffected by this
overall convention.
\end{remark}

\section{Relation to the XXZ chain and special limits}
\label{sec:eight-vertex-special-limits}

The XYZ chain is the most general nearest-neighbour spin-$1/2$ chain with
diagonal exchange couplings and no magnetic field:
\begin{equation}
\label{eq:XYZ-general-H}
    H_{\rm XYZ}
    =
    \sum_j
    \left(
        J_x\sigma_j^x\sigma_{j+1}^x
        +J_y\sigma_j^y\sigma_{j+1}^y
        +J_z\sigma_j^z\sigma_{j+1}^z
    \right).
\end{equation}
The XXZ chain is the special case
\begin{equation}
\label{eq:XXZ-as-XYZ-special-case}
    J_x=J_y,
    \qquad
    J_z=\Delta J_x .
\end{equation}
In the vertex-model language this corresponds to the degeneration
\begin{equation}
\label{eq:eight-to-six-d-zero}
    d(u)\to0,
\end{equation}
so that the eight-vertex $R$-matrix becomes a six-vertex $R$-matrix and the
elliptic parametrization degenerates to a trigonometric one.  In theta-function
language this occurs in the trigonometric limit of the elliptic nome.  In that
limit the $d$-weight disappears, $J_x-J_y\to0$, and the $U(1)$ symmetry of the
XXZ chain is restored.

\begin{center}
\begin{tabular}{@{}lll@{}}
\toprule
Vertex model & Quantum chain & Symmetry \\
\midrule
Six-vertex & XXZ & $U(1)$ generated by $S^z_{\rm tot}$ \\
Eight-vertex & XYZ & $\ZZ_2$ generated by $\prod_j\sigma_j^z$ \\
Trigonometric $R$ & XXZ anisotropy & $J_x=J_y$ \\
Elliptic $R$ & XYZ anisotropy & $J_x\neq J_y\neq J_z$ generically \\
\bottomrule
\end{tabular}
\end{center}

The transverse-field Ising and XY chains studied earlier in these notes are also
anisotropic spin chains, but their solvability mechanism is different.  They
become quadratic fermion models after the Jordan--Wigner transformation.  The
generic XYZ chain does not become a free-fermion Hamiltonian under
Jordan--Wigner; it becomes an interacting fermion model.  Its exact solvability
comes instead from the commuting transfer matrices inherited from the elliptic
Yang--Baxter equation.

\section{Fermionic image and loss of free-fermion structure}
\label{sec:XYZ-JW-image}

It is useful to compare the XYZ chain with the XY mother model.  Write
\begin{equation}
\label{eq:XYZ-in-ladder-operators}
\begin{aligned}
    J_x\sigma_j^x\sigma_{j+1}^x+J_y\sigma_j^y\sigma_{j+1}^y
    &=(J_x+J_y)
    \left(\sigma_j^+\sigma_{j+1}^-+\sigma_j^-\sigma_{j+1}^+\right)\\
    &\quad
    +(J_x-J_y)
    \left(\sigma_j^+\sigma_{j+1}^++\sigma_j^-\sigma_{j+1}^-\right).
\end{aligned}
\end{equation}
The first line gives fermion hopping after Jordan--Wigner, while the second line
gives fermion pairing.  The additional $J_z$ term gives
\begin{equation}
\label{eq:ZZ-to-density-interaction}
    \sigma_j^z\sigma_{j+1}^z
    =(1-2n_j)(1-2n_{j+1})
    =1-2(n_j+n_{j+1})+4n_jn_{j+1},
\end{equation}
up to the convention for which spin state is identified with an occupied
fermion.  Therefore the XYZ chain maps schematically to
\begin{equation}
\label{eq:XYZ-JW-schematic}
    H_{\rm XYZ}
    \longmapsto
    H_{\rm hopping}+H_{\rm pairing}+4J_z\sum_j n_jn_{j+1}
    +\text{chemical-potential and boundary terms}.
\end{equation}
The density-density interaction is absent only when $J_z=0$.  Hence the generic
XYZ chain is not a quadratic BdG problem.

\begin{remark}[Two different meanings of anisotropy]
In the free-fermion XY chain, anisotropy between $x$ and $y$ exchange produces
pairing but does not destroy quadratic solvability.  In the XYZ chain, the
$J_z\sigma^z\sigma^z$ term produces a genuine fermion interaction.  Therefore
``anisotropic spin chain'' does not by itself specify the solution method.
\end{remark}

\section{Baxter functional relations}
\label{sec:eight-vertex-functional-relations}

The six-vertex model can be diagonalized by coordinate or algebraic Bethe ansatz
because the transfer matrix has a simple reference state, such as the all-up
state.  The eight-vertex transfer matrix does not preserve the number of down
spins, so that reference-state construction is not directly available.
Baxter's solution instead uses functional relations.

The central object is a second commuting family of operators $Q(u)$ satisfying a
relation of the schematic form
\begin{equation}
\label{eq:eight-vertex-TQ-schematic}
    \tau(u)Q(u)
    =
    \phi_+(u)Q(u-2\eta)+\phi_-(u)Q(u+2\eta),
\end{equation}
where $\phi_{\pm}(u)$ are known scalar functions determined by the elliptic
weights and the system size.  On a common eigenvector of $\tau(u)$ and $Q(u)$,
this becomes a scalar finite-difference equation for the transfer-matrix
eigenvalue.

\begin{principle}[Baxter's replacement for the Bethe ansatz]
For the eight-vertex model, the role of Bethe roots is played by zeros of the
$Q$-eigenvalue.  The Bethe equations arise from requiring the apparent poles or
zeros in the $TQ$ relation to be consistent.  Thus the solution is still an
exact spectral solution, but it is naturally formulated as an elliptic functional
relation rather than as a magnon scattering problem.
\end{principle}

This is why the eight-vertex and XYZ models are often regarded as more advanced
than the six-vertex and XXZ models.  They are not merely more anisotropic; they
belong to the elliptic level of Yang--Baxter integrability.

\section{Summary of the model hierarchy}
\label{sec:eight-vertex-summary-hierarchy}

The hierarchy developed in the last two chapters can be summarized as follows:
\begin{equation}
\label{eq:R-matrix-hierarchy-summary}
\begin{array}{ccccc}
\text{XXX chain}
&\longleftrightarrow&
\text{rational }R\text{-matrix}
&\longleftrightarrow&
\Delta=1,\\[3pt]
\text{XXZ chain}
&\longleftrightarrow&
\text{trigonometric }R\text{-matrix}
&\longleftrightarrow&
J_x=J_y\neq J_z,\\[3pt]
\text{XYZ chain}
&\longleftrightarrow&
\text{elliptic }R\text{-matrix}
&\longleftrightarrow&
J_x\neq J_y\neq J_z.
\end{array}
\end{equation}
The same progression appears on the classical two-dimensional side:
\begin{equation}
\label{eq:vertex-model-hierarchy-summary}
    \text{six-vertex model}
    \subset
    \text{eight-vertex model},
\end{equation}
where the inclusion is obtained by setting the pair-flip weight $d$ to zero.

\boxedpar{%
The eight-vertex model is the elliptic vertex model whose regular transfer
matrix generates the XYZ chain.  Its solvability is not free-fermionic and not
based on $U(1)$ magnon-number conservation.  It is Yang--Baxter integrability at
the elliptic level, realized through commuting transfer matrices and Baxter
functional relations.}
